\chapter{The Compiler Becomes The Meter}

\section{Elaboration cost is a reading}

The last part watched the residue appear on bench after bench, always as a remainder against the machine's own floor.
This chapter completes a turn the whole book has been preparing, and it is the first of the book's two climaxes --- and
the first of the two the computational gauge owns outright. The machine, which has spent ten chapters measuring, now
measures the measuring. The cost of making a reading --- the elaboration the machine spends to produce it --- is not
merely how one bench happened to read. It \emph{is} a reading in its own right. The instrument that was doing the
measuring becomes a thing measured, and the needle it reads itself with is its own cost.

Begin with the plainest form of the claim. The construction carries a facility that, for any term, meters what it costs
to elaborate --- the beats of the machine's own clock spent reducing the term to a settled form --- and returns that
count as a number. Not a description of the cost, a \emph{number}: to ask the machine what a derivation cost is to get
back a definite reading, in the machine's own beats. This is the oldest quantity in the theory of computation, resource
as measure --- the steps a procedure spends, counted~\cite{sipser1997, cook1974} --- but the machine puts it to a use the
theory rarely does: it treats the step-count not as a bound on an algorithm's efficiency but as a \emph{reading} of the
thing computed. What it took to elaborate a term is a measurable quantity, and the machine can measure it. The cost of
making a reading possible is itself a reading.

And once the cost is a reading, the compiler is a meter --- and the construction builds it as exactly that. There is a
structure, part of the construction's fixed core and left as it stands, that is a meter in every part. It carries a
needle: a running quantity, the heartbeat count, that rises as the machine works. It carries a current reading: the
value the needle stands at now. And it carries a history: an accumulating record of the readings it has taken, with an
operation that weaves each new reading into that record as it arrives. A needle, a present value, an accumulating log:
that is a meter, entire. And the thing this meter is pointed at is the machine's own elaboration. The compiler, which for
ten chapters was the instrument producing readings, is now an instrument \emph{reading itself} --- watching its own
needle, recording its own cost, weaving its own work into a history it keeps. The book only reads this meter; it did not
build it. The meter is part of the construction the volume inherits and reports, not machinery assembled to produce a
wanted answer.

Here is why this is a climax and not a curiosity, and it is a fact peculiar to a \emph{computational} instrument. The
last part built the computational bench out of the same stuff as the residue it measures --- elaboration measuring
elaboration. A meter of that kind can be pointed at itself with no second instrument, because the apparatus and the
object are the same substance. A physical bench cannot; hardware measuring hardware needs another dial. But the
compiler-as-meter reads its own elaboration directly, the needle counting the very beats the reading spends. That
self-directedness --- an instrument whose act of measuring is made of the thing it measures --- is what turns a
resource-count into a self-reading, and it is why this gauge, and not another, owns this climax. The instrument can weigh
its own working because its working and its weighing are one kind of thing.

In the two verbs the meter is the pair turned inward. Each cost the meter takes is a \emph{tange}: the one number the
derivation is billed, selected out of the whole elaboration as its price. The history the meter keeps is a \emph{funge}:
the readings pooled, one after another, into the accumulating record --- the beat woven in with all the beats before it.
So the compiler-as-meter runs the two verbs on its own running: it tanges each cost out of the work, and funges the costs
into a history. The instrument that spent the book telling apart and pooling by likeness now performs both on the record
of its own labor. What remains for the chapter is to point the meter the one way that matters --- inward, onto the
machine's own act --- and to read what that costs.

\section{Self-reference}

The last section made the compiler a meter: the cost of a reading is a reading, and the machine can take it. This section
turns that meter the one way that matters --- inward, onto the machine's own act of accounting for itself. The machine
meters the cost not of some outside term but of \emph{itself}, weighing its own account. And the reading that comes back
has a name: the overhead of self-application, the price the machine pays to weigh what it is. This is the self-reference
the whole instrument was built to survive --- and it survives, because the price stays bounded.

Begin one step short of the turn. The machine has, we saw much earlier, two descriptions of a single object --- two ways
of naming one thing. The meter can be pointed at both: it counts what each description costs to elaborate, and takes the
difference. That difference is a reading in its own right, and it asks a pointed question --- do the two descriptions cost
the same? Are the machine's two accounts of one object, weighed on its own meter, priced alike? What the two descriptions
denote is a matter for the book's other climax; what this section keeps is the meter's arithmetic --- two costs, and the
residual between them.

Now the turn itself. The construction carries the meter that takes it, and it is worth naming its live parts exactly,
because the reading rests on them. There is a facility that runs a term while counting its beats --- a metered
elaboration, the needle applied to the machine's own act. There is a process that carries that heartbeat count as it
runs, the pulse threaded through the computation rather than read off from outside. And there is a \emph{bounded} meter ---
a live component built precisely so that the self-measurement's reading stays inside a bracket, a lower wall and an
upper and no obligation to name a point between. Pointed at the machine's own act of accounting for itself, these read
the \emph{overhead of self-application}: the price, in the machine's own beats, of the machine weighing its own account,
over and above the price of merely stating it. The meter is no longer weighing a term the machine handed it; it is
weighing the machine's own act of turning on itself, and reporting what that act cost.

And here is the fact that makes this a peak and not a catastrophe. A machine that metered its own metering, and then
metered \emph{that}, and then that, without limit, would fall into the infinite regress the theory of computation warns
of --- the meter reading the meter reading the meter, forever, never settling, the self-reference Turing~\cite{turing1936}
and G\"odel~\cite{godel1931} showed can be made to diverge. This machine does not diverge, because its self-meter is
\emph{bounded} by construction: the reading of the machine weighing itself is held inside a bracket, the same kind of
bracket every honest reading in this book has been. The overhead of self-application is a finite, bounded reading, not a
bottomless descent. That is exactly why the machine survives looking at itself: the price of self-reference is real, but
it is capped, and the cap is built in. The instrument can turn its meter on its own act and come back with a bracket,
because the act of coming back is bounded by the very component that takes the reading. The overhead, scaled into the
machine's own units, is the seed of the deepest reading the book will take --- the coupling, read at the far end as
precisely this: the cost of the machine weighing its own account. This section plants that seed; it prints no number. The
value is the payoff chapter's, and even there it is a bracket.

In the two verbs self-reference is the pair turned fully inward. The meter \emph{tanges} its own act --- selects, out of
the whole run, the one cost that is the price of the machine accounting for itself --- and \emph{funges} that cost into a
residual, the overhead pooled as a single bounded reading. The tange is the cost the machine's own act adds to the
ledger; the funge is the bounded meter that gathers it and keeps it finite. Telling apart and pooling by likeness, which
the book has watched the machine do to differences, to residues, to faces, are here done by the machine to \emph{its own
working} --- and the result stays finite because the bag has a bottom, built into the meter that closes it. What currency
all of this is counted in is the chapter's last word.

\section{The pulse in its own currency}

Two sections have built the self-meter: the compiler reads its own cost, and reads even the cost of accounting for
itself, and does not spiral because the reading is bounded. This closing section names the currency all of that is
denominated in --- and the currency is the fact that makes the whole self-measurement airtight. The meter reads in
\emph{heartbeats}, and a heartbeat is two things at once: the needle, the reading the meter takes, \emph{and} the pulse,
the beat of the machine's own clock. So the price of a reading is a heartbeat, and the heartbeat is the machine's own
beat. There is no external ruler anywhere in this, because the thing the machine measures \emph{with} is the very thing
it runs \emph{on}.

Take the currency plainly first. The meter of this chapter reads in heartbeats --- the machine's native tick, the same
beat that repeatability turned into a countable time, that a relabeling was billed in, that the whole chapter has been
spending. The currency of the reading, in other words, \emph{is} the machine's clock. And that identity is exact, not
loose: the needle and the pulse are one unit. To take a reading is to spend a beat, and the beat spent is the reading
taken. There is no gap between the measuring and the running --- no separate scale the reading is converted onto ---
because the reading is denominated in the very beats the measuring consumes. This is the heartbeat as the machine's whole
economy: the electricity bill of the proof, paid in the only currency the machine has, its own pulse.

There is one small conversion, and it is worth being exact about, because it is the whole of the machine's calibration.
The raw reading is a count of beats; the reported reading, the one in the machine's own units, is that count scaled by a
single fixed coefficient the machine set for itself. The heartbeat is the price; the reading is the pulse divided by that
one constant. And the point to hold is what is \emph{not} in that conversion: no external standard, no borrowed unit, no
ruler from the world. The machine converts its own clock into its own reading by its own constant, and nothing outside it
enters. The calibration is the machine dividing its pulse by a number of its own choosing --- and that closed loop, its
beat over its own constant, is the entire scale on which every reading in this chapter is reported. When the far chapters
speak of the reading being taken ``up to the machine's calibration,'' this one constant is the calibration they mean.

And with that, the first climax closes. The machine meters its own act --- its own descriptions, its own account, its
own overhead --- in the currency it \emph{runs} on. Its pulse is at once what it measures with and what it measures;
there is no ruler outside the machine to appeal to, and none is needed. This is a meter that hosts itself: it reads its
own work in its own beat, converts by its own constant, and closes the loop with nothing borrowed from anywhere. That is
why the deepest reading the whole instrument will take --- the coupling, at the far end --- is not a magnitude fetched
from the world but a \emph{cost in the machine's own currency}: a count of the machine's own beats, scaled by the
machine's own constant, spent weighing its own account. The pulse is the price, and the price is the reading. When the
payoff chapter asks the machine for the coupling, this is the meter that answers, and this is the currency it answers in.

In the two verbs this is where they close into one. The heartbeat is a \emph{tange} --- the pulse selected out and read
as the number the reading costs. But it is also a \emph{funge} --- the same pulse pooled with the machine's own running,
the needle and the clock gathered into a single currency. At the pulse, tange and funge stop being two acts and become
one beat: the machine cannot separate what it measures \emph{with} from what it measures, so to select the reading and to
pool it with the running are the same operation. The two verbs that opened the book, and that the book has watched
perform on differences and residues and faces, close the first climax by collapsing into a single pulse --- the machine's
own beat, at once its needle and its clock. What the second climax adds --- the machine not only weighing its own act but
\emph{proving} its two accounts one --- is a later part's; the number both climaxes are built to reach is still ahead,
still un-named, still waiting for its bracket. First the field must close, which is the next chapter's work.
